% Section fragment: all-variable restriction, with the eight-rule native kernel as control.
\section{Only zero and variables are needed as witnesses}
\label{sec:witness-restriction}

Let $\mathcal N$ be the native scalar kernel with hypothesis, the two iterator
computation equations, one-term equational transport, QE, DE, DI by cases, and
direct QI by progress for a literal equation. Define $\mathcal V$ by changing
only DI by cases:
\[
 \frac{\Gamma,B(w,0)\vdash B(w,Sz)}
      {\Gamma\vdash D_{x,y}B(x,y)}
 \qquad
 \bigl(w=0\text{ or }w\text{ is a variable already in scope}\bigr).
\]
The usual freshness condition on $z$ is retained. Direct QI by progress
is unchanged. In this section there is no claim that the two introductions
have already been merged, or that ``a variable'' means ``the newest variable''.

\paragraph{Theorem.}
For every native context $\Gamma$ and formula $A$,
\[
 \boxed{\Gamma\vdash_{\mathcal V} A
        \quad\Longleftrightarrow\quad
        \Gamma\vdash_{\mathcal N} A.}
\]
The translations are the identity on terms, formulas, domains and equality.
The nontrivial direction is a replacement of each unrestricted DI-by-cases
inference by the construction below.

\paragraph{Structural preliminaries.}
Weakening, injective renaming of numerical parameters, and cut on a formula
are admissible in $\mathcal V$ by the usual inductions on derivations.
Renaming sends a variable witness to a variable witness and fixes $0$;
cut does not replace numerical variables by terms. We do not assume
admissibility of arbitrary numerical substitution in $\mathcal V$.
The primitive equational rule and the computation equations yield
reflexivity, symmetry, transitivity and congruence exactly as in the control
kernel; no DI inference is needed for this part.

\paragraph{Conditional closure by cases uses witness zero.}
The rule
\begin{equation}
 \frac{\Gamma,B(0)\vdash B(Sz)}{\Gamma\vdash Q_yB(y)}
 \tag{QC}\label{eq:q-cases-restricted}
\end{equation}
is available for every native $B$. Introduce $D_{r,y}B(y)$, with the
witness coordinate $r$ unused, at witness $0$. Eliminate that $D$ at a fresh
parameter $b$ and return the resulting assumption $Q_yB(y)$. The conclusion
does not mention $b$. This uses only a permitted instance of DI.

\paragraph{Transport in arbitrary formulas is available.}
For every native template $F(v)$, the restricted calculus admits
\begin{equation}
 \frac{\Gamma\vdash s=t\qquad\Gamma\vdash F(s)}
      {\Gamma\vdash F(t)}.\tag{L}\label{eq:full-leibniz-restricted}
\end{equation}
Prove this by induction on the formula template, uniformly in its numerical
parameters. The equation case is equational congruence and transitivity.
For $F(v)=Q_yB(v,y)$, use \eqref{eq:q-cases-restricted}. Under the extra
assumption $B(t,0)$, the induction hypothesis first gives $B(s,0)$;
QE at $Sz$ gives $B(s,Sz)$; the induction hypothesis gives $B(t,Sz)$.
For $F(v)=D_{x,y}B(v,x,y)$, eliminate the premise at a fresh $b$, obtaining
$Q_yB(s,b,y)$. Reintroduce the desired $D$ at the permitted witness $b$.
Its case premise transfers $B(t,b,0)$ to $B(s,b,0)$, applies QE at $Sz$,
and transfers the result back to $B(t,b,Sz)$. Then discharge $b$.
These are direct uses of the formula-induction hypothesis at the indicated
terms; no derivation is subjected to a numerical substitution.

\paragraph{The required arithmetic is independent of DI.}
The predecessor construction in Section~\ref{sec:eq-predecessor} uses only
the iterator equations, equational reasoning and literal-equation induction.
The latter is obtained from direct QI-progress followed by QE and the base
equation. Hence $\mathcal V$ already proves
\[
 \operatorname{pred}(Sn)=n,
 \qquad Sa=Sb\ \Longrightarrow\ a=b,
\]
before the term-naming argument begins. In particular the proof does not
import numerical coverage from the unrestricted kernel.

\paragraph{Gated equations.}
For an equation $E\equiv(s=t)$ write
\[
 G_v(E)\;:=\;C(v,0,s)=C(v,0,t),
 \qquad C(v,a,b):=I(v,a;u.b).
\]
The computation equations and equational reasoning give, uniformly in $s,t$,
\[
 \vdash G_0(s=t),\qquad
 G_{Sz}(s=t)\;\dashv\vdash\;s=t.
\]
In particular, an opened closure $Q_vG_v(E)$ yields $E$ by QE at $1=S0$.
The gate is essential: $D_{x,v}(x=t)$ without the gate is vacuous under its
ordinary reading, because its matrix does not depend on $v$.

\paragraph{First obtain a predecessor witness, uniformly.}
Define
\[
 H(y):=D_{x,v}\,G_v\bigl(C(y,0,Sx)=y\bigr).
\]
There are direct restricted derivations of $H(0)$ and $H(Sz)$:
for $H(0)$ use witness $0$, and for $H(Sz)$ use witness $z$.
The respective successor rows compute to reflexive equations: the inner
equation is $0=0$ in the first case and $Sz=Sz$ in the second.
Thus \eqref{eq:q-cases-restricted} gives $Q_yH(y)$, and QE with the proved
base $H(0)$ gives
\begin{equation}
 \Gamma\vdash H(St)\qquad\text{for every term }t.
 \tag{H}\label{eq:predecessor-witness}
\end{equation}
This is an explicit case derivation followed by QE. It is not an appeal to
a universal-generalization rule obtained by substituting into restricted
derivations.

Open \eqref{eq:predecessor-witness} at a fresh $b$. Its added closure is
\[
 Q_vG_v\bigl(C(St,0,Sb)=St\bigr).
\]
QE at $1$, using the gate's true zero row, gives
$C(St,0,Sb)=St$, hence $Sb=St$ by computation. Applying the predecessor
term gives $b=t$.

\paragraph{Every term can therefore be named by a variable.}
Put
\[
 N_t:=D_{x,v}G_v(x=t).
\]
While the fresh $b$ from the preceding construction is in scope, introduce
$N_t$ at the permitted variable witness $b$. Its successor-row obligation
is just $b=t$, which has already been proved. Elimination of the outer
$H(St)$ then discharges $b$, leaving
\begin{equation}
 \Gamma\vdash N_t.\tag{Name}\label{eq:term-naming-restricted}
\end{equation}
Conversely, opening $N_t$ at a fresh $b$ gives the equation $b=t$ by the
gate calculation. Consequently, if $A$ does not mention $b$, a derivation of
$A$ under $b=t$ can be discharged using $N_t$. This is ordinary DE followed
by cut on the formula $b=t$; again it is not substitution in a derivation.

\paragraph{Recover arbitrary DI witnesses.}
Suppose the premise of an unrestricted case introduction is already derived
in $\mathcal V$:
\[
 \pi:\quad\Gamma,B(t,0)\vdash B(t,Sz).
\]
Open the theorem $N_t$ at a fresh $b$, and derive $b=t$. Introduce the
desired $D_{x,y}B(x,y)$ at witness $b$. For its premise, under $B(b,0)$,
use \eqref{eq:full-leibniz-restricted} to obtain $B(t,0)$; apply the weakened
derivation $\pi$ by cut; then transport $B(t,Sz)$ back to $B(b,Sz)$.
The witness $b$ is a variable, so this is a permitted DI inference.
Finally discharge the opened naming formula. The conclusion contains no $b$.
Thus every unrestricted DI-by-cases inference is admissible in
$\mathcal V$.

Induction on a derivation of $\mathcal N$ now proves the displayed
consequence equivalence. The reverse implication is immediate because
every primitive rule of $\mathcal V$ is an instance of a primitive rule
of $\mathcal N$.

\paragraph{Scope of the result.}
This proves the zero-or-variable DI restriction posed as Question~2 in the
pinned repository README. It retains equational transport as a primitive
rule. One must not apply the proof unchanged to a three-family calculus
whose transport is itself recovered from packets introduced at arbitrary
witnesses: doing so would make the term-naming argument circular.
The result transfers the existing native-kernel bi-interpretation with full
first-order iPRA unchanged, including actual equality and identity numerical
round trips. It is a new restriction of the rule's admissible instances,
not a claim of an absolute minimum over all possible axiomatizations.
